\section{Introdução}\label{sec:introducao}

A madeira acompanha a história da construção desde os primeiros sistemas estruturais e permanece atual em razão de sua ecoeficiência, do baixo consumo de energia em sua produção e da elevada relação resistência-peso \parencite{NEVES2025, FRAGA2025}. Ao longo do tempo, o material passou a ser combinado a outros sistemas, como as estruturas mistas de madeira-concreto e madeira-aço, ampliando as possibilidades de racionalização de seu emprego \parencite{MIOTTO2009}. Diante das metas globais de descarbonização da construção civil, a madeira ressurge como alternativa estrutural promissora para pontes de pequeno e médio porte, particularmente em estradas vicinais, nas quais os vãos são reduzidos e a demanda por soluções econômicas e sustentáveis é elevada.

O uso da madeira como sistema estrutural confiável decorre da aplicação criteriosa de normas técnicas, que estabelecem os parâmetros para o dimensionamento e a verificação da segurança \parencite{ADOLFS2011}. No Brasil, a ABNT NBR 7190 \parencite{NBR7190-1:2022} regulamenta o projeto de estruturas de madeira e fornece as diretrizes para o cálculo das seções e a verificação dos estados limites. Contudo, o dimensionamento desses elementos envolve múltiplas variáveis de projeto, tais como dimensões geométricas, propriedades do material e condições de carregamento, o que torna o procedimento complexo e fortemente dependente da experiência do projetista.

Historicamente, o processo de projeto estrutural evoluiu de uma abordagem clássica, determinística, baseada em cálculos manuais e em procedimentos semiautomatizados de desenho assistido por computador, para uma abordagem denominada projeto inteligente, que integra modelagem paramétrica, análise estrutural automatizada e técnicas de otimização orientadas a dados \parencite{DEB2009}. Nesse novo paradigma, o dimensionamento deixa de ser um procedimento linear e passa a ser tratado como um processo iterativo e integrado, no qual o modelo geométrico, a análise dos esforços e a decisão de projeto são acoplados em uma única plataforma computacional. Essa transição permite reduzir o tempo de projeto, ampliar a qualidade das soluções e explorar de forma sistemática o espaço de alternativas viáveis.

Nesse contexto, o dimensionamento de pontes de madeira pode ser formulado como um problema de otimização, no qual objetivos conflitantes devem ser considerados simultaneamente, como a minimização do consumo de material, o atendimento aos estados limites último e de serviço e o desempenho estrutural global. A otimização multiobjetivo baseada em metaheurísticas surge como ferramenta adequada para lidar com essa complexidade, pois possibilita a obtenção de um conjunto de soluções eficientes representadas por uma fronteira de Pareto, oferecendo ao projetista maior flexibilidade na tomada de decisão e favorecendo o uso racional dos recursos \parencite{blank2020}.

Entretanto, a otimização determinística conduz, com frequência, a soluções localizadas exatamente sobre a fronteira das restrições, nas quais pequenas variações nas variáveis de projeto, decorrentes de tolerâncias de fabricação, variabilidade dimensional da madeira roliça ou imprecisões de montagem, podem levar à violação dos critérios de segurança. Esse aspecto é particularmente crítico em madeira roliça, cujo diâmetro apresenta variabilidade natural significativa. Assim, torna-se necessário incorporar o conceito de robustez ao processo de otimização, de modo que as soluções obtidas permaneçam seguras mesmo diante de desvios em torno de seus valores nominais.

Este trabalho propõe uma abordagem de projeto inteligente para o dimensionamento de pontes de madeira roliça, na qual a modelagem paramétrica dos elementos estruturais é acoplada a uma rotina de otimização robusta multiobjetivo, em conformidade com os critérios da ABNT NBR 7190 \parencite{NBR7190-1:2022}. A robustez é incorporada de forma nativa ao processo de busca, avaliando cada solução candidata sob um desvio de 5\% aplicado aos valores nominais das variáveis de projeto. A metodologia é implementada em uma plataforma computacional de acesso livre, capaz de automatizar o processo de dimensionamento, gerar a fronteira eficiente e emitir relatórios completos de verificação.

Os objetivos específicos do trabalho são: (i) formular o dimensionamento dos elementos estruturais de pontes de madeira como um problema de otimização multiobjetivo, definindo funções objetivo e restrições compatíveis com os critérios de segurança e desempenho da ABNT NBR 7190 \parencite{NBR7190-1:2022}; (ii) implementar o algoritmo NSGA-II acoplado ao modelo estrutural paramétrico, assegurando a convergência e a diversidade das soluções; (iii) incorporar a análise de robustez, avaliando o efeito de um desvio de 5\% nas variáveis de projeto sobre as funções objetivo e as restrições; e (iv) demonstrar a aplicação da metodologia em pontes de estrada vicinal submetidas a diferentes trens tipo, comparando o desempenho da otimização robusta com o de um dimensionamento sem estratégia de otimização.
